\begin{frame}{복습 주간 일정 (시험 전 1주)}
  \small
  \begin{tabular}{@{}clp{0.62\textwidth}@{}}
    \toprule
    요일 & 시간 & 할 일 \\
    \midrule
    월 & 30분 & 20문항 자가진단 1차 $\to$ ✗가 모인 2$\sim$3개 장 특정 \\
    화 & 2시간 & §유도 템플릿 3가지의 예제를 \textbf{전부 손으로} 풀고
      노트북으로 검증 \\
    수 & 2시간 & 특정 장의 ``자주 하는 실수'' + 해당 장 확인 문제 재독 \\
    목 & 90분 & 8.1 프로젝트 노트북을 \textbf{처음부터 재실행} ---
      ``코드의 각 단계를 설명할 수 있는가''에 집중 \\
    금 & 1시간 & 이 절 확인 문제 4개 + 8.1 확인 문제 4개를
      \textbf{참고 없이} 풀어보기 \\
    토 오전 & 30분 & 20문항 자가진단 2차 --- 월요일과 ✗ 개수 비교 \\
    일 & --- & \textbf{새로운 내용 금지} --- 충분한 수면 \\
    \bottomrule
  \end{tabular}
  \vskip0.6em
  \begin{block}{순서가 왜 이 모양인가}
    \textbf{유도가 재독보다 먼저} --- 손으로 유도한 뒤에야 ``막힌
    지점''이 명확해져 재독이 ``다 읽어보기''가 아니라 ``막힌 지점만
    파기''가 됨. \textbf{실험이 마지막} --- 프로젝트 코드가 ``새로운
    MDP에 개념을 적용'' 그 자체의 가장 좋은 연습이고, 공식이 손에
    익은 뒤에야 코드 각 줄에 붙을 개념이 보이기 때문.
  \end{block}
\end{frame}
